\begin{helper}
Let \(S\) be a finite set and let \(p,x\in\mathbb R\).  Then
\[
\sum_{B\subseteq S}p^{|B|}(1-p)^{|S|-|B|}x^{|B|}
=((1-p)+px)^{|S|}.
\]
\end{helper}
\begin{proof}
For each subset \(B\subseteq S\), combine the two factors depending on
\(|B|\):
\[
p^{|B|}x^{|B|}=(px)^{|B|}.
\]
Therefore the displayed sum is
\[
\sum_{B\subseteq S}(px)^{|B|}(1-p)^{|S|-|B|}.
\]
The finite binomial theorem, written in subset form, says that for any real
numbers \(r\) and \(q\),
\[
\sum_{B\subseteq S}r^{|B|}q^{|S|-|B|}=(r+q)^{|S|}.
\]
Applying this with \(r=px\) and \(q=1-p\) gives
\[
\sum_{B\subseteq S}(px)^{|B|}(1-p)^{|S|-|B|}
=((px)+(1-p))^{|S|}.
\]
Finally, addition is commutative, so \((px)+(1-p)=(1-p)+px\), which is the
claimed identity.
\end{proof}
